%!TEX root = ../../../main.tex

\chapter{Teststrategie und Evaluation}
\label{chap:evaluation}

Wenn man mit einem Agenten chattet, wirkt er schnell gut, weil die Antworten flüssig klingen. Ob die Zahl darin stimmt, sieht man nicht. Ich teste deshalb nur mit festen Fragen, jede mit Typ, Sollwert, Solldokument und Begründung. Die Läufe sind wiederholbar (Temperatur 0, fester Seed, festes \enquote{heutiges} Datum für Fristen). Es gibt drei Arten von Testsätzen. Der \emph{Kernsatz} enthält die geforderten fünf Standardfragen, zwei Grenzfälle und zwei Negativfragen (Tabelle~\ref{tab:kern}). Am \emph{Entwicklungssatz} (54 Fragen) suche ich Fehler. \emph{Versiegelte Sätze} mit 20 bis 24 neuen Fragen schreibe ich vor einem größeren Umbau, sichere sie per SHA-256 und messe sie nur einmal. Nötig wurde das nach einer unangenehmen Messung. Auf dem Entwicklungssatz kam der Orchestrator auf 81\,\%, auf dem ersten versiegelten Satz nur auf 35\,\%, weil ich meine Regeln unbewusst auf meine eigenen Fragen zugeschnitten hatte. Die Ergebnisse unten stammen vom letzten versiegelten Satz mit 24 Fragen (Tabelle~\ref{tab:testsatz}), der Orchestrator lief zweimal, der ReAct-Agent wegen der Laufzeit einmal.

\begin{table}[tb]
    \centering\small
    \begin{tabular}{@{}l r r@{}}
        \toprule
        & \textbf{Orchestrator} & \textbf{ReAct-Agent} \\
        \midrule
        Korrektheit gesamt & 77\,\% (37/48) & 33\,\% (8/24) \\
        \quad 95-\%-Intervall (Wilson) & 63--87\,\% & 18--53\,\% \\
        \quad a / b / c (Wert, Text, Vergleich) & 27/30 & 3/15 \\
        \quad d (Listen) & 2/10 & 1/5 \\
        \quad e (ohne Beleg, richtig abgelehnt) & 8/8 & 4/4 \\
        Antworten mit richtiger Quelle & 36/40 & 7/13 \\
        Antwortzeit Median / 95.~Perzentil & 98 / 199\,s & 92 / 318\,s \\
        \midrule
        RAGAS Faithfulness (20 / 13 Antworten) & 0,68 [0,51--0,84] & 0,60 [0,35--0,83] \\
        RAGAS Context Recall (15 Fragen, a--c) & 0,93 [0,80--1,00] & 0,20 [0,00--0,40] \\
        \bottomrule
    \end{tabular}
    \caption{Ergebnisse auf dem letzten versiegelten Satz (24 Fragen, Orchestrator zwei Läufe; RAGAS mit 95-\%-Bootstrap-Intervall)}
    \label{tab:ergebnisse}
\end{table}

\section{Automatisierte RAGAS-Evaluation}
\label{sec:ragas}

Eine RAG-Antwort scheitert entweder an der Suche oder am Lesen. RAGAS bewertet beides getrennt. Die Metriken sind probabilistisch, denn ein Sprachmodell urteilt, bei mir Qwen2.5 (7B).

\paragraph{Faithfulness.} Der Richter zerlegt die Antwort in einzelne Aussagen (\enquote{Der Grundpreis beträgt 12,50\,€ und entspricht dem Vertrag} sind zwei) und prüft für jede, ob sie aus den benutzten Belegen folgt. Faithfulness ist der Anteil dieser belegten Aussagen an allen Aussagen, 1,0 heißt, dass nichts erfunden ist. Für mich ist das die wichtigste Zahl, weil eine überzeugende, aber falsche Aussage über Geld schlimmer ist als keine Antwort.

\paragraph{Context Recall.} Hier zerlegt der Richter die Sollantwort, die ich aus der Belegstelle bilde (\enquote{Laut Dokument: 2,80\,€/m³}), und misst, welcher Anteil ihrer Sätze in den gefundenen Belegen steht. Ist eine Antwort falsch, obwohl der Recall 1 ist, hat das Modell falsch gelesen. Den Recall gibt es nur für die Typen a bis c, Listen und Negativfragen haben keine einzelne Belegstelle.

\paragraph{Ergebnisse.} Der Context Recall des Orchestrators liegt bei 0,93 (Tabelle~\ref{tab:ergebnisse}). Die Suche findet also fast immer die richtige Stelle. Beim ReAct-Agenten sind es nur 0,20, er wählt seine Suchen selbst und kommt oft gar nicht dort an. Die Faithfulness von 0,68 kam mir zuerst sehr niedrig vor, beim Nachlesen der Bewertungen hat sich das Bild aber verschoben. Ein Teil der Abzüge ist berechtigt. Bei T13 schreibt der Orchestrator zum Beispiel nur \enquote{entspricht dem Preis im Vertrag}, ohne den Preis von 33,60\,€ zu nennen. Ein anderer Teil liegt am Richter. Die richtige und wörtlich belegte Antwort \enquote{251,37\,€} auf T01 bekam bei beiden Metriken 0. Bei T04 hatte ich aus zwei Läufen fast gleiche Antworten. \enquote{\dots{} brutto mit 294,25\,€ berechnet [1]} bekam 1,0, \enquote{\dots{} mit 294,25\,€ berechnet, wie in Beleg [1] ersichtlich} bekam 0,0.

\paragraph{RAGAS gegen meine Zahlprüfung.} Ich habe den Richter mit einer deterministischen Prüfung verglichen, die jede Zahl meldet, die in keinem Beleg vorkommt. Bei den 20 Antworten schlugen beide zweimal gemeinsam an, beide Male bei Listen. Achtmal schlug nur RAGAS an, und vier dieser acht Antworten waren laut Sollwert richtig. Zweimal schlug nur die Zahlprüfung an. Bei T07 hatte das Modell \enquote{1. Januar 2024} ausgeschrieben, die Prüfung kannte aber nur 01.01.2024, und bei T14 hatte das Modell selbst gerechnet. Ein 7B-Richter zeigt also, wo man hinschauen sollte, als alleinige Note taugt er nicht.

\section{Validierung mittels Golden Dataset}
\label{sec:golden}

Beim Golden Dataset bewertet kein Modell, sondern Code. Eine Antwort ist richtig, wenn jeder Sollwert in einem Satz steht, der ihn nicht verneint. Bei Listen müssen alle erwarteten Einträge dabei sein, bei Negativfragen zählt nur eine Ablehnung. Weil 24 Fragen wenig sind, gebe ich ein 95-\%-Intervall nach Wilson an. Eine erfundene Begründung neben der richtigen Zahl bemerkt die Regel nicht, dafür gibt es Faithfulness.

\paragraph{Tool- und Modultests.} Darunter liegen 130 automatische Tests vom Parser bis zur Zahlprüfung. Für den Planer laufen nach jeder Änderung gespeicherte Pläne von 106 alten Fragen ohne Modellaufruf durch die Code-Prüfung, und keine darf fälschlich abgelehnt werden. Das hat zum Beispiel meine erste Regel gegen die Haftpflicht-Falle gestoppt, die nebenbei auch \enquote{Wie viel Umsatzsteuer \dots?} abgelehnt hätte.

\paragraph{Ergebnisse.} Auf dem Kernsatz sind acht von neun Antworten richtig. Bei K7 listet der Orchestrator zusätzlich Verträge ohne Angabe zur Verlängerung auf. Auf dem versiegelten Satz liegt der Orchestrator in drei von vier Fällen richtig, der ReAct-Agent nur in jedem dritten. Fragen ohne passendes Dokument lehnen beide immer ab, auch Fallen wie \enquote{Versicherung für den VW Golf} (das Auto steht nur in einer Werkstattrechnung). Die klare Schwäche sind Listen mit 2 von 10. Der Python-Filter rechnet richtig, das Problem liegt davor im Plan. Bei \enquote{Rechnungen zwischen 100 und 1\,000\,€} übernimmt das Modell nur eine der beiden Grenzen, bei \enquote{SEPA-Verträge von Max Mustermann} nur eine der beiden Bedingungen. Mein Planer erlaubt nur eine Bedingung je Filter.

\paragraph{Was ich daraus mitnehme.} Wo der Code entscheidet, also bei Fristen, Filtern und der Bestandsprüfung, hat das System fast immer gestimmt. Fehler entstehen, wo das Modell liest oder einen Filter beschreibt, und der ReAct-Agent, bei dem das Modell alles entscheidet, liegt deutlich darunter. Für ein 7--8B-Modell auf der CPU halte ich einen festen Ablauf mit Prüfungen im Code deshalb für die bessere Wahl. Allerdings stammen Dokumente und Fragen von mir, die Testsätze sind klein, und echte Unterlagen sind oft gescannt oder unvollständig. Für meinen eigenen Ordner bräuchte LocalLLM noch Texterkennung, Filter mit mehreren Bedingungen und eine verschlüsselte Ablage.
